\section*{अध्याय \ref{ch:g2:where-animals-live} --- जंतु कहाँ रहते हैं}

\begin{solution}{exo:g2:where-animals-live:1}
वह तरह की जगह जहाँ \omterm{def:g1:animals-around-us:animal}{जंतु} स्वाभाविक रूप से रहता है। वहाँ उसे भोजन, पानी,
ओट, और अपने बच्चों के लिए एक सुरक्षित कोना मिलता है।
\end{solution}

\begin{solution}{exo:g2:where-animals-live:2}
तालाब: मेंढक, व्याध-पतंगा (या काँटा-मछली, तालाब का घोंघा)। जंगल:
गिलहरी, कठफोड़वा (या हिरन, जंगली चींटी)। मैदान: ख़रगोश, टिड्डा (या
तितली, खेत का चूहा)।
\end{solution}

\begin{solution}{exo:g2:where-animals-live:3}
काँटा-मछली: तालाब। कठफोड़वा: जंगल। टिड्डा: मैदान। कठ-जूँ: पुरानी दीवार
के पत्थरों के नीचे (नम, छिपे हुए कोने)।
\end{solution}

\begin{solution}{exo:g2:where-animals-live:4}
मेंढक: झिल्लीदार पैर, जो पानी को धकेलते हैं। गिलहरी: पकड़ने वाले पंजे
और संतुलन बनाती बड़ी पूँछ।
\end{solution}

\begin{solution}{exo:g2:where-animals-live:5}
चौड़े फावड़े जैसे खोदने वाले हाथ; दोनों दिशाओं में चपटे लेटते \omterm{def:g1:animals-around-us:covering}{रोम},
ताकि मिट्टी उन्हें कभी न उलझाए; और वे \omterm{def:g1:the-five-senses:sense}{आँखें} जिनकी अँधेरे में उसे
ज़रूरत ही नहीं।
\end{solution}

\begin{solution}{exo:g2:where-animals-live:6}
कठफोड़वा: तनों पर ऊँचाई पर। गिलहरी: शाखाओं में। हिरन: झाड़ियों के बीच।
जंगली चींटी: ज़मीन पर, अपने ढेर जैसे घर में।
\end{solution}

\begin{solution}{exo:g2:where-animals-live:7}
वे तालाब के साथ ही ग़ायब हो जाते हैं --- \omterm{def:g2:where-animals-live:habitat}{आवास} के बिना वे भोजन, ओट और
प्रजनन की जगह, तीनों एक साथ खो देते हैं। इसलिए \omterm{def:g1:animals-around-us:animal}{जंतुओं} की रक्षा उनकी
रहने की जगहों की रक्षा से शुरू होती है।
\end{solution}

\begin{solution}{exo:g2:where-animals-live:8}
उत्तर खुला है। एक ही उलटे पत्थर के नीचे अकसर कठ-जूएँ, चींटियाँ, कोई
भृंग और कभी-कभी कनखजूरा मिलते हैं --- छोटे \omterm{def:g2:where-animals-live:habitat}{आवास} हैरान कर देने वाली
हद तक भरे होते हैं।
\end{solution}

\begin{solution}{exo:g2:where-animals-live:9}
झिल्लीदार पैर (खेने के लिए), पानी बहा देने वाले \omterm{def:g1:animals-around-us:covering}{पंख}, छानने वाली चोंच:
तीनों पानी की ओर इशारा करते हैं --- तालाब या झील का \omterm{def:g2:where-animals-live:habitat}{आवास}। \omterm{def:g1:animals-around-us:animal}{जंतु} बहुत
संभव है कि बत्तख़ हो।
\end{solution}

\begin{solution}{exo:g2:where-animals-live:10}
मेंढक का साज़-सामान --- झिल्लीदार पैर, नम \omterm{def:g1:the-five-senses:sense}{त्वचा}, पानी और नम किनारों
के लिए बना शरीर --- सूखे डिब्बे में बेकार है, और उसका \omterm{def:g2:where-animals-live:habitat}{आवास} उसे भोजन
से कहीं ज़्यादा देता था: उसकी \omterm{def:g1:the-five-senses:sense}{त्वचा} और उसके \omterm{def:g2:animals-grow-up:born}{अंडों} के लिए पानी, छिपने
की जगहें, दूसरे मेंढक। किसी जंगली \omterm{def:g1:animals-around-us:animal}{जंतु} पर दया यही है कि उसे उसके \omterm{def:g2:where-animals-live:habitat}{आवास}
में रहने दिया जाए --- या वहाँ लौटा दिया जाए।
\end{solution}
